% ============================================================================
% (c) 2025 Mohamed EL AFRIT -- IPSSI
% Licence CC BY-NC-SA 4.0
% https://creativecommons.org/licenses/by-nc-sa/4.0/deed.fr
%
% FICHE DE SYNTHESE -- Jour 4 : Reseaux de neurones
% ============================================================================
% ============================================================================
% © 2025 Mohamed EL AFRIT — IPSSI
% Ce contenu est distribué sous licence Creative Commons BY-NC-SA 4.0
% https://creativecommons.org/licenses/by-nc-sa/4.0/deed.fr
% ============================================================================
%
% TEMPLATE COMMUN — Mathématiques appliquées à l'Intelligence Artificielle
% Ce fichier est le préambule partagé de TOUS les documents LaTeX du projet.
%
% Utilisation :
%   % ============================================================================
% © 2025 Mohamed EL AFRIT — IPSSI
% Ce contenu est distribué sous licence Creative Commons BY-NC-SA 4.0
% https://creativecommons.org/licenses/by-nc-sa/4.0/deed.fr
% ============================================================================
%
% TEMPLATE COMMUN — Mathématiques appliquées à l'Intelligence Artificielle
% Ce fichier est le préambule partagé de TOUS les documents LaTeX du projet.
%
% Utilisation :
%   % ============================================================================
% © 2025 Mohamed EL AFRIT — IPSSI
% Ce contenu est distribué sous licence Creative Commons BY-NC-SA 4.0
% https://creativecommons.org/licenses/by-nc-sa/4.0/deed.fr
% ============================================================================
%
% TEMPLATE COMMUN — Mathématiques appliquées à l'Intelligence Artificielle
% Ce fichier est le préambule partagé de TOUS les documents LaTeX du projet.
%
% Utilisation :
%   \input{template_ipssi}          % Importer ce préambule
%   \jourNumero{1}                  % Définir le numéro du jour (1, 2, 3 ou 4)
%   \jourTheme{Données et représentation mathématique}
%   \begin{document}
%   \pagegarde{Fiche de synthèse}{1}{Données et représentation mathématique}
%   ...
%   \end{document}
% ============================================================================

% --- Classe de document ---
\documentclass[11pt,a4paper]{article}

% ============================================================================
% ENCODAGE ET LANGUE
% ============================================================================
\usepackage[utf8]{inputenc}
\usepackage[T1]{fontenc}
\usepackage[french]{babel}

% ============================================================================
% MATHÉMATIQUES
% ============================================================================
\usepackage{amsmath,amssymb,amsfonts,mathtools}

% Commandes mathématiques personnalisées (conventions du cours)
\newcommand{\vx}{\mathbf{x}}          % vecteur x
\newcommand{\vw}{\mathbf{w}}          % vecteur w (poids)
\newcommand{\vy}{\mathbf{y}}          % vecteur y
\newcommand{\mX}{\mathbf{X}}          % matrice X
\newcommand{\mW}{\mathbf{W}}          % matrice W
\newcommand{\yhat}{\hat{y}}           % prédiction
\newcommand{\pscal}[2]{\langle #1, #2 \rangle} % produit scalaire
\newcommand{\gradJ}{\nabla J(\boldsymbol{\theta})} % gradient de J

% ============================================================================
% MISE EN PAGE
% ============================================================================
\usepackage[margin=2.5cm]{geometry}
\usepackage{setspace}
\setstretch{1.15}                      % interligne légèrement aéré

% ============================================================================
% COULEURS
% ============================================================================
\usepackage[dvipsnames,svgnames,x11names]{xcolor}

% Palette principale du projet
\definecolor{ipssiBleu}{HTML}{3B82F6}       % Définitions
\definecolor{ipssiVert}{HTML}{22C55E}       % À retenir
\definecolor{ipssiOrange}{HTML}{F97316}     % Attention / Piège
\definecolor{ipssiViolet}{HTML}{A855F7}     % Intuition
\definecolor{ipssiCyan}{HTML}{06B6D4}       % Exemples
\definecolor{ipssiGrisClair}{HTML}{F3F4F6}  % Fond code
\definecolor{ipssiGrisFonce}{HTML}{1F2937}  % Texte code
\definecolor{ipssiFondPage}{HTML}{FAFAFA}   % Fond léger page de garde
\definecolor{ipssiRouge}{HTML}{E74C3C}      % Classe 0 (salaire ≤ 45k€)
\definecolor{ipssiVertData}{HTML}{2ECC71}   % Classe 1 (salaire > 45k€)
\definecolor{ipssiBleuDroite}{HTML}{3498DB} % Droite de régression

% ============================================================================
% ENCADRÉS PÉDAGOGIQUES (tcolorbox)
% ============================================================================
\usepackage[most]{tcolorbox}

% Style de base partagé
\tcbset{
    ipssibase/.style={
        enhanced,
        breakable,
        left=4mm,
        right=4mm,
        top=3mm,
        bottom=3mm,
        boxrule=0pt,
        frame hidden,
        borderline west={3pt}{0pt}{#1},
        colback=#1!5,
        colframe=#1,
        fonttitle=\bfseries\sffamily,
        coltitle=#1!80!black,
        attach boxed title to top left={yshift=-2mm, xshift=4mm},
        boxed title style={
            colback=#1!15,
            colframe=#1,
            boxrule=0.5pt,
            arc=2pt,
        },
        arc=2pt,
        outer arc=2pt,
    }
}

% Icônes TikZ pour les encadrés (compatible pdflatex, pas d'emoji Unicode)
\newcommand{\icoDefinition}{%
    \tikz[baseline=-0.3em]{\draw[ipssiBleu, thick] (0,0) -- (0.3,0) -- (0.3,0.35) -- (0,0.35) -- cycle;
    \draw[ipssiBleu, thick] (0.05,0.1) -- (0.25,0.1); \draw[ipssiBleu, thick] (0.05,0.2) -- (0.2,0.2);}%
}
\newcommand{\icoRetenir}{%
    \tikz[baseline=-0.2em]{\draw[ipssiVert, thick, fill=ipssiVert!20] (0.15,0) -- (0.3,0.3) -- (0,0.3) -- cycle;}%
}
\newcommand{\icoAttention}{%
    \tikz[baseline=-0.2em]{\draw[ipssiOrange, thick] (0.15,0.35) -- (0,0) -- (0.3,0) -- cycle;
    \node[ipssiOrange, font=\tiny\bfseries] at (0.15,0.12) {!};}%
}
\newcommand{\icoIntuition}{%
    \tikz[baseline=-0.2em]{\draw[ipssiViolet, thick, fill=ipssiViolet!15]
    (0.15,0.35) -- (0.08,0.2) -- (0.08,0.15) -- (0.22,0.15) -- (0.22,0.2) -- cycle;
    \draw[ipssiViolet, thick] (0.11,0.12) -- (0.11,0.05); \draw[ipssiViolet, thick] (0.19,0.12) -- (0.19,0.05);}%
}
\newcommand{\icoExemple}{%
    \tikz[baseline=-0.2em]{\draw[ipssiCyan, thick] (0,0) rectangle (0.3,0.3);
    \draw[ipssiCyan, thick] (0,0.15) -- (0.3,0.15); \draw[ipssiCyan, thick] (0.15,0) -- (0.15,0.3);}%
}
\newcommand{\icoPython}{%
    \tikz[baseline=-0.2em]{\node[font=\small\bfseries\ttfamily, ipssiVert!70!black] at (0,0) {\textgreater\textgreater};}%
}

% Définition (bleu)
\newtcolorbox{definitionbox}[1][D\'{e}finition]{
    ipssibase=ipssiBleu,
    title={\icoDefinition\; #1},
}

% À retenir (vert)
\newtcolorbox{retenirbox}[1][\`{A} retenir]{
    ipssibase=ipssiVert,
    title={\icoRetenir\; #1},
}

% Attention / Piège (orange)
\newtcolorbox{attentionbox}[1][Attention]{
    ipssibase=ipssiOrange,
    title={\icoAttention\; #1},
}

% Intuition (violet)
\newtcolorbox{intuitionbox}[1][Intuition]{
    ipssibase=ipssiViolet,
    title={\icoIntuition\; #1},
}

% Exemple (cyan)
\newtcolorbox{exemplebox}[1][Exemple]{
    ipssibase=ipssiCyan,
    title={\icoExemple\; #1},
}

% Code Python (style terminal)
\newtcolorbox{pythonbox}[1][Code Python]{
    enhanced,
    breakable,
    left=4mm, right=4mm, top=3mm, bottom=3mm,
    boxrule=0.5pt,
    colback=ipssiGrisClair,
    colframe=ipssiGrisFonce!40,
    fonttitle=\bfseries\ttfamily\small,
    coltitle=ipssiGrisFonce,
    title={\icoPython\; #1},
    attach boxed title to top left={yshift=-2mm, xshift=4mm},
    boxed title style={colback=ipssiGrisClair, colframe=ipssiGrisFonce!40, boxrule=0.5pt, arc=2pt},
    arc=2pt, outer arc=2pt,
}

% ============================================================================
% CODE PYTHON (listings)
% ============================================================================
\usepackage{listings}

\lstdefinestyle{python}{
    language=Python,
    basicstyle=\ttfamily\small,
    keywordstyle=\color{ipssiBleu}\bfseries,
    stringstyle=\color{ipssiVert!70!black},
    commentstyle=\color{gray}\itshape,
    numberstyle=\tiny\color{gray},
    numbers=left,
    numbersep=8pt,
    backgroundcolor=\color{ipssiGrisClair},
    frame=none,
    rulecolor=\color{ipssiGrisFonce!20},
    breaklines=true,
    breakatwhitespace=true,
    tabsize=4,
    showstringspaces=false,
    captionpos=b,
    aboveskip=8pt,
    belowskip=8pt,
    xleftmargin=10pt,
    framexleftmargin=10pt,
    morekeywords={as, with, True, False, None, self, np, plt, import, from},
    literate=
        {é}{{\'e}}1 {è}{{\`e}}1 {ê}{{\^e}}1 {ë}{{\"e}}1
        {à}{{\`a}}1 {â}{{\^a}}1
        {ù}{{\`u}}1 {û}{{\^u}}1
        {ô}{{\^o}}1
        {î}{{\^i}}1 {ï}{{\"i}}1
        {ç}{{\c{c}}}1
        {É}{{\'E}}1 {È}{{\`E}}1
        {→}{{$\rightarrow$}}1
        {≈}{{$\approx$}}1
        {≤}{{$\leq$}}1
        {≥}{{$\geq$}}1
        {★}{{$\bigstar$}}1
        {ŷ}{{$\hat{y}$}}1
        {·}{{$\cdot$}}1,
}

\lstset{style=python}

% Style pour le code inline
\newcommand{\pyinline}[1]{\lstinline[style=python,basicstyle=\ttfamily\small]{#1}}

% ============================================================================
% GRAPHIQUES
% ============================================================================
\usepackage{tikz}
\usepackage{pgfplots}
\pgfplotsset{compat=1.18}
\usetikzlibrary{arrows.meta, calc, positioning, shapes.geometric, decorations.pathreplacing}

% ============================================================================
% TABLEAUX
% ============================================================================
\usepackage{booktabs}
\usepackage{tabularx}
\usepackage{multirow}

% ============================================================================
% LISTES
% ============================================================================
\usepackage{enumitem}
\setlist[itemize]{itemsep=2pt, parsep=0pt, topsep=4pt}
\setlist[enumerate]{itemsep=2pt, parsep=0pt, topsep=4pt}

% ============================================================================
% HYPERLIENS
% ============================================================================
\usepackage[
    colorlinks=true,
    linkcolor=ipssiBleu!80!black,
    urlcolor=ipssiBleu!80!black,
    citecolor=ipssiVert!80!black,
    bookmarks=true,
    bookmarksopen=true,
    pdfauthor={Mohamed EL AFRIT},
    pdftitle={Mathématiques appliquées à l'IA — IPSSI},
    pdfsubject={Formation Bachelor 2 — IPSSI 2025-2026},
]{hyperref}

% ============================================================================
% IMAGES
% ============================================================================
\usepackage{graphicx}

% ============================================================================
% VARIABLES PARAMÉTRABLES (jour et thème)
% ============================================================================
\newcommand{\leJour}{X}
\newcommand{\leTheme}{Thème}
\newcommand{\jourNumero}[1]{\renewcommand{\leJour}{#1}}
\newcommand{\jourTheme}[1]{\renewcommand{\leTheme}{#1}}

% ============================================================================
% EN-TÊTE ET PIED DE PAGE (fancyhdr)
% ============================================================================
\usepackage{fancyhdr}

\pagestyle{fancy}
\fancyhf{}

% En-tête
\fancyhead[L]{\small\sffamily\textcolor{gray}{IPSSI — Maths appliquées à l'IA}}
\fancyhead[R]{\small\sffamily\textcolor{gray}{Jour \leJour\ — \leTheme}}
\renewcommand{\headrulewidth}{0.4pt}

% Pied de page
\fancyfoot[L]{\footnotesize\textcolor{gray}{Mohamed EL AFRIT\enspace\textbullet\enspace\url{www.mohamedelafrit.com}\enspace\textbullet\enspace CC BY-NC-SA 4.0}}
\fancyfoot[C]{\footnotesize\textcolor{gray}{\thepage}}
\fancyfoot[R]{\footnotesize\textcolor{gray}{2025–2026}}
\renewcommand{\footrulewidth}{0.2pt}

% Style pour la page de garde (pas d'en-tête)
\fancypagestyle{pagegarde}{
    \fancyhf{}
    \renewcommand{\headrulewidth}{0pt}
    \fancyfoot[C]{\footnotesize\textcolor{gray}{\thepage}}
    \renewcommand{\footrulewidth}{0pt}
}

% ============================================================================
% PAGE DE GARDE — \pagegarde{titre}{jour}{theme}
% ============================================================================
\newcommand{\pagegarde}[3]{%
    \thispagestyle{pagegarde}
    \begin{center}
        \vspace*{1.5cm}

        % Logo / Bandeau institution
        {\Large\sffamily\bfseries\textcolor{ipssiBleu}{IPSSI}}\par
        \vspace{2pt}
        {\normalsize\sffamily 2\textsuperscript{e} année Bachelor Informatique}\par

        \vspace{0.8cm}
        {\color{ipssiBleu}\rule{0.6\textwidth}{1.5pt}}\par
        \vspace{0.8cm}

        % Titre du module
        {\LARGE\sffamily\bfseries Mathématiques appliquées\\[4pt] à l'Intelligence Artificielle}\par

        \vspace{1cm}

        % Titre du document
        \begin{tcolorbox}[
            enhanced,
            colback=ipssiBleu!5,
            colframe=ipssiBleu!50,
            boxrule=1pt,
            arc=4pt,
            width=0.85\textwidth,
            halign=center,
            top=10pt, bottom=10pt,
        ]
            {\huge\sffamily\bfseries\textcolor{ipssiBleu!80!black}{#1}}\par
            \vspace{6pt}
            {\Large\sffamily Jour #2 — #3}
        \end{tcolorbox}

        \vspace{1.5cm}

        % Informations auteur
        {\large\sffamily
            \textbf{Auteur :} Mohamed EL AFRIT\\[4pt]
            \url{www.mohamedelafrit.com}\\[8pt]
            \textbf{Année académique :} 2025–2026
        }\par

        \vspace{1.5cm}

        % Licence Creative Commons
        {\color{ipssiBleu}\rule{0.4\textwidth}{0.5pt}}\par
        \vspace{8pt}
        {\small\sffamily
            \textcopyright\ 2025 Mohamed EL AFRIT — IPSSI\\[2pt]
            Ce contenu est distribué sous licence\\[2pt]
            \textbf{Creative Commons BY-NC-SA 4.0}\\[2pt]
            \url{https://creativecommons.org/licenses/by-nc-sa/4.0/deed.fr}
        }\par
    \end{center}
    \newpage
}

% ============================================================================
% NIVEAUX D'EXERCICES
% ============================================================================

% Étoile compatible pdflatex
\newcommand{\starfull}[1]{\textcolor{#1}{$\bigstar$}}
\newcommand{\starempty}{\textcolor{gray!40}{$\bigstar$}}

% Fondamental (1 étoile)
\newcommand{\exofondamental}{%
    \starfull{ipssiVert}\starempty\starempty\enspace%
    {\small\sffamily\textcolor{ipssiVert}{\textbf{Fondamental}}}%
}

% Intermédiaire (2 étoiles)
\newcommand{\exointermediaire}{%
    \starfull{ipssiOrange}\starfull{ipssiOrange}\starempty\enspace%
    {\small\sffamily\textcolor{ipssiOrange}{\textbf{Interm\'{e}diaire}}}%
}

% Avancé (3 étoiles)
\newcommand{\exoavance}{%
    \starfull{ipssiRouge}\starfull{ipssiRouge}\starfull{ipssiRouge}\enspace%
    {\small\sffamily\textcolor{ipssiRouge}{\textbf{Avanc\'{e}}}}%
}

% Commande d'exercice numéroté
\newcounter{exocount}[section]
\newcommand{\exercice}[1]{%
    \stepcounter{exocount}%
    \vspace{8pt}%
    \noindent\textbf{\sffamily Exercice \theexocount}\enspace — \enspace #1\par
    \vspace{4pt}%
}

% ============================================================================
% COMMANDES UTILITAIRES
% ============================================================================

% Séparateur visuel entre sections
\newcommand{\separateur}{%
    \vspace{6pt}
    \begin{center}
        {\color{ipssiBleu!30}\rule{0.3\textwidth}{0.5pt}}
    \end{center}
    \vspace{6pt}
}

% Mot-clé en gras coloré
\newcommand{\motcle}[1]{\textbf{\textcolor{ipssiBleu!80!black}{#1}}}

% Formule encadrée importante
\newcommand{\formule}[1]{%
    \begin{center}
        \fcolorbox{ipssiBleu!50}{ipssiBleu!5}{%
            \quad$\displaystyle #1$\quad%
        }
    \end{center}
}

% ============================================================================
% FIN DU PRÉAMBULE
% ============================================================================
          % Importer ce préambule
%   \jourNumero{1}                  % Définir le numéro du jour (1, 2, 3 ou 4)
%   \jourTheme{Données et représentation mathématique}
%   \begin{document}
%   \pagegarde{Fiche de synthèse}{1}{Données et représentation mathématique}
%   ...
%   \end{document}
% ============================================================================

% --- Classe de document ---
\documentclass[11pt,a4paper]{article}

% ============================================================================
% ENCODAGE ET LANGUE
% ============================================================================
\usepackage[utf8]{inputenc}
\usepackage[T1]{fontenc}
\usepackage[french]{babel}

% ============================================================================
% MATHÉMATIQUES
% ============================================================================
\usepackage{amsmath,amssymb,amsfonts,mathtools}

% Commandes mathématiques personnalisées (conventions du cours)
\newcommand{\vx}{\mathbf{x}}          % vecteur x
\newcommand{\vw}{\mathbf{w}}          % vecteur w (poids)
\newcommand{\vy}{\mathbf{y}}          % vecteur y
\newcommand{\mX}{\mathbf{X}}          % matrice X
\newcommand{\mW}{\mathbf{W}}          % matrice W
\newcommand{\yhat}{\hat{y}}           % prédiction
\newcommand{\pscal}[2]{\langle #1, #2 \rangle} % produit scalaire
\newcommand{\gradJ}{\nabla J(\boldsymbol{\theta})} % gradient de J

% ============================================================================
% MISE EN PAGE
% ============================================================================
\usepackage[margin=2.5cm]{geometry}
\usepackage{setspace}
\setstretch{1.15}                      % interligne légèrement aéré

% ============================================================================
% COULEURS
% ============================================================================
\usepackage[dvipsnames,svgnames,x11names]{xcolor}

% Palette principale du projet
\definecolor{ipssiBleu}{HTML}{3B82F6}       % Définitions
\definecolor{ipssiVert}{HTML}{22C55E}       % À retenir
\definecolor{ipssiOrange}{HTML}{F97316}     % Attention / Piège
\definecolor{ipssiViolet}{HTML}{A855F7}     % Intuition
\definecolor{ipssiCyan}{HTML}{06B6D4}       % Exemples
\definecolor{ipssiGrisClair}{HTML}{F3F4F6}  % Fond code
\definecolor{ipssiGrisFonce}{HTML}{1F2937}  % Texte code
\definecolor{ipssiFondPage}{HTML}{FAFAFA}   % Fond léger page de garde
\definecolor{ipssiRouge}{HTML}{E74C3C}      % Classe 0 (salaire ≤ 45k€)
\definecolor{ipssiVertData}{HTML}{2ECC71}   % Classe 1 (salaire > 45k€)
\definecolor{ipssiBleuDroite}{HTML}{3498DB} % Droite de régression

% ============================================================================
% ENCADRÉS PÉDAGOGIQUES (tcolorbox)
% ============================================================================
\usepackage[most]{tcolorbox}

% Style de base partagé
\tcbset{
    ipssibase/.style={
        enhanced,
        breakable,
        left=4mm,
        right=4mm,
        top=3mm,
        bottom=3mm,
        boxrule=0pt,
        frame hidden,
        borderline west={3pt}{0pt}{#1},
        colback=#1!5,
        colframe=#1,
        fonttitle=\bfseries\sffamily,
        coltitle=#1!80!black,
        attach boxed title to top left={yshift=-2mm, xshift=4mm},
        boxed title style={
            colback=#1!15,
            colframe=#1,
            boxrule=0.5pt,
            arc=2pt,
        },
        arc=2pt,
        outer arc=2pt,
    }
}

% Icônes TikZ pour les encadrés (compatible pdflatex, pas d'emoji Unicode)
\newcommand{\icoDefinition}{%
    \tikz[baseline=-0.3em]{\draw[ipssiBleu, thick] (0,0) -- (0.3,0) -- (0.3,0.35) -- (0,0.35) -- cycle;
    \draw[ipssiBleu, thick] (0.05,0.1) -- (0.25,0.1); \draw[ipssiBleu, thick] (0.05,0.2) -- (0.2,0.2);}%
}
\newcommand{\icoRetenir}{%
    \tikz[baseline=-0.2em]{\draw[ipssiVert, thick, fill=ipssiVert!20] (0.15,0) -- (0.3,0.3) -- (0,0.3) -- cycle;}%
}
\newcommand{\icoAttention}{%
    \tikz[baseline=-0.2em]{\draw[ipssiOrange, thick] (0.15,0.35) -- (0,0) -- (0.3,0) -- cycle;
    \node[ipssiOrange, font=\tiny\bfseries] at (0.15,0.12) {!};}%
}
\newcommand{\icoIntuition}{%
    \tikz[baseline=-0.2em]{\draw[ipssiViolet, thick, fill=ipssiViolet!15]
    (0.15,0.35) -- (0.08,0.2) -- (0.08,0.15) -- (0.22,0.15) -- (0.22,0.2) -- cycle;
    \draw[ipssiViolet, thick] (0.11,0.12) -- (0.11,0.05); \draw[ipssiViolet, thick] (0.19,0.12) -- (0.19,0.05);}%
}
\newcommand{\icoExemple}{%
    \tikz[baseline=-0.2em]{\draw[ipssiCyan, thick] (0,0) rectangle (0.3,0.3);
    \draw[ipssiCyan, thick] (0,0.15) -- (0.3,0.15); \draw[ipssiCyan, thick] (0.15,0) -- (0.15,0.3);}%
}
\newcommand{\icoPython}{%
    \tikz[baseline=-0.2em]{\node[font=\small\bfseries\ttfamily, ipssiVert!70!black] at (0,0) {\textgreater\textgreater};}%
}

% Définition (bleu)
\newtcolorbox{definitionbox}[1][D\'{e}finition]{
    ipssibase=ipssiBleu,
    title={\icoDefinition\; #1},
}

% À retenir (vert)
\newtcolorbox{retenirbox}[1][\`{A} retenir]{
    ipssibase=ipssiVert,
    title={\icoRetenir\; #1},
}

% Attention / Piège (orange)
\newtcolorbox{attentionbox}[1][Attention]{
    ipssibase=ipssiOrange,
    title={\icoAttention\; #1},
}

% Intuition (violet)
\newtcolorbox{intuitionbox}[1][Intuition]{
    ipssibase=ipssiViolet,
    title={\icoIntuition\; #1},
}

% Exemple (cyan)
\newtcolorbox{exemplebox}[1][Exemple]{
    ipssibase=ipssiCyan,
    title={\icoExemple\; #1},
}

% Code Python (style terminal)
\newtcolorbox{pythonbox}[1][Code Python]{
    enhanced,
    breakable,
    left=4mm, right=4mm, top=3mm, bottom=3mm,
    boxrule=0.5pt,
    colback=ipssiGrisClair,
    colframe=ipssiGrisFonce!40,
    fonttitle=\bfseries\ttfamily\small,
    coltitle=ipssiGrisFonce,
    title={\icoPython\; #1},
    attach boxed title to top left={yshift=-2mm, xshift=4mm},
    boxed title style={colback=ipssiGrisClair, colframe=ipssiGrisFonce!40, boxrule=0.5pt, arc=2pt},
    arc=2pt, outer arc=2pt,
}

% ============================================================================
% CODE PYTHON (listings)
% ============================================================================
\usepackage{listings}

\lstdefinestyle{python}{
    language=Python,
    basicstyle=\ttfamily\small,
    keywordstyle=\color{ipssiBleu}\bfseries,
    stringstyle=\color{ipssiVert!70!black},
    commentstyle=\color{gray}\itshape,
    numberstyle=\tiny\color{gray},
    numbers=left,
    numbersep=8pt,
    backgroundcolor=\color{ipssiGrisClair},
    frame=none,
    rulecolor=\color{ipssiGrisFonce!20},
    breaklines=true,
    breakatwhitespace=true,
    tabsize=4,
    showstringspaces=false,
    captionpos=b,
    aboveskip=8pt,
    belowskip=8pt,
    xleftmargin=10pt,
    framexleftmargin=10pt,
    morekeywords={as, with, True, False, None, self, np, plt, import, from},
    literate=
        {é}{{\'e}}1 {è}{{\`e}}1 {ê}{{\^e}}1 {ë}{{\"e}}1
        {à}{{\`a}}1 {â}{{\^a}}1
        {ù}{{\`u}}1 {û}{{\^u}}1
        {ô}{{\^o}}1
        {î}{{\^i}}1 {ï}{{\"i}}1
        {ç}{{\c{c}}}1
        {É}{{\'E}}1 {È}{{\`E}}1
        {→}{{$\rightarrow$}}1
        {≈}{{$\approx$}}1
        {≤}{{$\leq$}}1
        {≥}{{$\geq$}}1
        {★}{{$\bigstar$}}1
        {ŷ}{{$\hat{y}$}}1
        {·}{{$\cdot$}}1,
}

\lstset{style=python}

% Style pour le code inline
\newcommand{\pyinline}[1]{\lstinline[style=python,basicstyle=\ttfamily\small]{#1}}

% ============================================================================
% GRAPHIQUES
% ============================================================================
\usepackage{tikz}
\usepackage{pgfplots}
\pgfplotsset{compat=1.18}
\usetikzlibrary{arrows.meta, calc, positioning, shapes.geometric, decorations.pathreplacing}

% ============================================================================
% TABLEAUX
% ============================================================================
\usepackage{booktabs}
\usepackage{tabularx}
\usepackage{multirow}

% ============================================================================
% LISTES
% ============================================================================
\usepackage{enumitem}
\setlist[itemize]{itemsep=2pt, parsep=0pt, topsep=4pt}
\setlist[enumerate]{itemsep=2pt, parsep=0pt, topsep=4pt}

% ============================================================================
% HYPERLIENS
% ============================================================================
\usepackage[
    colorlinks=true,
    linkcolor=ipssiBleu!80!black,
    urlcolor=ipssiBleu!80!black,
    citecolor=ipssiVert!80!black,
    bookmarks=true,
    bookmarksopen=true,
    pdfauthor={Mohamed EL AFRIT},
    pdftitle={Mathématiques appliquées à l'IA — IPSSI},
    pdfsubject={Formation Bachelor 2 — IPSSI 2025-2026},
]{hyperref}

% ============================================================================
% IMAGES
% ============================================================================
\usepackage{graphicx}

% ============================================================================
% VARIABLES PARAMÉTRABLES (jour et thème)
% ============================================================================
\newcommand{\leJour}{X}
\newcommand{\leTheme}{Thème}
\newcommand{\jourNumero}[1]{\renewcommand{\leJour}{#1}}
\newcommand{\jourTheme}[1]{\renewcommand{\leTheme}{#1}}

% ============================================================================
% EN-TÊTE ET PIED DE PAGE (fancyhdr)
% ============================================================================
\usepackage{fancyhdr}

\pagestyle{fancy}
\fancyhf{}

% En-tête
\fancyhead[L]{\small\sffamily\textcolor{gray}{IPSSI — Maths appliquées à l'IA}}
\fancyhead[R]{\small\sffamily\textcolor{gray}{Jour \leJour\ — \leTheme}}
\renewcommand{\headrulewidth}{0.4pt}

% Pied de page
\fancyfoot[L]{\footnotesize\textcolor{gray}{Mohamed EL AFRIT\enspace\textbullet\enspace\url{www.mohamedelafrit.com}\enspace\textbullet\enspace CC BY-NC-SA 4.0}}
\fancyfoot[C]{\footnotesize\textcolor{gray}{\thepage}}
\fancyfoot[R]{\footnotesize\textcolor{gray}{2025–2026}}
\renewcommand{\footrulewidth}{0.2pt}

% Style pour la page de garde (pas d'en-tête)
\fancypagestyle{pagegarde}{
    \fancyhf{}
    \renewcommand{\headrulewidth}{0pt}
    \fancyfoot[C]{\footnotesize\textcolor{gray}{\thepage}}
    \renewcommand{\footrulewidth}{0pt}
}

% ============================================================================
% PAGE DE GARDE — \pagegarde{titre}{jour}{theme}
% ============================================================================
\newcommand{\pagegarde}[3]{%
    \thispagestyle{pagegarde}
    \begin{center}
        \vspace*{1.5cm}

        % Logo / Bandeau institution
        {\Large\sffamily\bfseries\textcolor{ipssiBleu}{IPSSI}}\par
        \vspace{2pt}
        {\normalsize\sffamily 2\textsuperscript{e} année Bachelor Informatique}\par

        \vspace{0.8cm}
        {\color{ipssiBleu}\rule{0.6\textwidth}{1.5pt}}\par
        \vspace{0.8cm}

        % Titre du module
        {\LARGE\sffamily\bfseries Mathématiques appliquées\\[4pt] à l'Intelligence Artificielle}\par

        \vspace{1cm}

        % Titre du document
        \begin{tcolorbox}[
            enhanced,
            colback=ipssiBleu!5,
            colframe=ipssiBleu!50,
            boxrule=1pt,
            arc=4pt,
            width=0.85\textwidth,
            halign=center,
            top=10pt, bottom=10pt,
        ]
            {\huge\sffamily\bfseries\textcolor{ipssiBleu!80!black}{#1}}\par
            \vspace{6pt}
            {\Large\sffamily Jour #2 — #3}
        \end{tcolorbox}

        \vspace{1.5cm}

        % Informations auteur
        {\large\sffamily
            \textbf{Auteur :} Mohamed EL AFRIT\\[4pt]
            \url{www.mohamedelafrit.com}\\[8pt]
            \textbf{Année académique :} 2025–2026
        }\par

        \vspace{1.5cm}

        % Licence Creative Commons
        {\color{ipssiBleu}\rule{0.4\textwidth}{0.5pt}}\par
        \vspace{8pt}
        {\small\sffamily
            \textcopyright\ 2025 Mohamed EL AFRIT — IPSSI\\[2pt]
            Ce contenu est distribué sous licence\\[2pt]
            \textbf{Creative Commons BY-NC-SA 4.0}\\[2pt]
            \url{https://creativecommons.org/licenses/by-nc-sa/4.0/deed.fr}
        }\par
    \end{center}
    \newpage
}

% ============================================================================
% NIVEAUX D'EXERCICES
% ============================================================================

% Étoile compatible pdflatex
\newcommand{\starfull}[1]{\textcolor{#1}{$\bigstar$}}
\newcommand{\starempty}{\textcolor{gray!40}{$\bigstar$}}

% Fondamental (1 étoile)
\newcommand{\exofondamental}{%
    \starfull{ipssiVert}\starempty\starempty\enspace%
    {\small\sffamily\textcolor{ipssiVert}{\textbf{Fondamental}}}%
}

% Intermédiaire (2 étoiles)
\newcommand{\exointermediaire}{%
    \starfull{ipssiOrange}\starfull{ipssiOrange}\starempty\enspace%
    {\small\sffamily\textcolor{ipssiOrange}{\textbf{Interm\'{e}diaire}}}%
}

% Avancé (3 étoiles)
\newcommand{\exoavance}{%
    \starfull{ipssiRouge}\starfull{ipssiRouge}\starfull{ipssiRouge}\enspace%
    {\small\sffamily\textcolor{ipssiRouge}{\textbf{Avanc\'{e}}}}%
}

% Commande d'exercice numéroté
\newcounter{exocount}[section]
\newcommand{\exercice}[1]{%
    \stepcounter{exocount}%
    \vspace{8pt}%
    \noindent\textbf{\sffamily Exercice \theexocount}\enspace — \enspace #1\par
    \vspace{4pt}%
}

% ============================================================================
% COMMANDES UTILITAIRES
% ============================================================================

% Séparateur visuel entre sections
\newcommand{\separateur}{%
    \vspace{6pt}
    \begin{center}
        {\color{ipssiBleu!30}\rule{0.3\textwidth}{0.5pt}}
    \end{center}
    \vspace{6pt}
}

% Mot-clé en gras coloré
\newcommand{\motcle}[1]{\textbf{\textcolor{ipssiBleu!80!black}{#1}}}

% Formule encadrée importante
\newcommand{\formule}[1]{%
    \begin{center}
        \fcolorbox{ipssiBleu!50}{ipssiBleu!5}{%
            \quad$\displaystyle #1$\quad%
        }
    \end{center}
}

% ============================================================================
% FIN DU PRÉAMBULE
% ============================================================================
          % Importer ce préambule
%   \jourNumero{1}                  % Définir le numéro du jour (1, 2, 3 ou 4)
%   \jourTheme{Données et représentation mathématique}
%   \begin{document}
%   \pagegarde{Fiche de synthèse}{1}{Données et représentation mathématique}
%   ...
%   \end{document}
% ============================================================================

% --- Classe de document ---
\documentclass[11pt,a4paper]{article}

% ============================================================================
% ENCODAGE ET LANGUE
% ============================================================================
\usepackage[utf8]{inputenc}
\usepackage[T1]{fontenc}
\usepackage[french]{babel}

% ============================================================================
% MATHÉMATIQUES
% ============================================================================
\usepackage{amsmath,amssymb,amsfonts,mathtools}

% Commandes mathématiques personnalisées (conventions du cours)
\newcommand{\vx}{\mathbf{x}}          % vecteur x
\newcommand{\vw}{\mathbf{w}}          % vecteur w (poids)
\newcommand{\vy}{\mathbf{y}}          % vecteur y
\newcommand{\mX}{\mathbf{X}}          % matrice X
\newcommand{\mW}{\mathbf{W}}          % matrice W
\newcommand{\yhat}{\hat{y}}           % prédiction
\newcommand{\pscal}[2]{\langle #1, #2 \rangle} % produit scalaire
\newcommand{\gradJ}{\nabla J(\boldsymbol{\theta})} % gradient de J

% ============================================================================
% MISE EN PAGE
% ============================================================================
\usepackage[margin=2.5cm]{geometry}
\usepackage{setspace}
\setstretch{1.15}                      % interligne légèrement aéré

% ============================================================================
% COULEURS
% ============================================================================
\usepackage[dvipsnames,svgnames,x11names]{xcolor}

% Palette principale du projet
\definecolor{ipssiBleu}{HTML}{3B82F6}       % Définitions
\definecolor{ipssiVert}{HTML}{22C55E}       % À retenir
\definecolor{ipssiOrange}{HTML}{F97316}     % Attention / Piège
\definecolor{ipssiViolet}{HTML}{A855F7}     % Intuition
\definecolor{ipssiCyan}{HTML}{06B6D4}       % Exemples
\definecolor{ipssiGrisClair}{HTML}{F3F4F6}  % Fond code
\definecolor{ipssiGrisFonce}{HTML}{1F2937}  % Texte code
\definecolor{ipssiFondPage}{HTML}{FAFAFA}   % Fond léger page de garde
\definecolor{ipssiRouge}{HTML}{E74C3C}      % Classe 0 (salaire ≤ 45k€)
\definecolor{ipssiVertData}{HTML}{2ECC71}   % Classe 1 (salaire > 45k€)
\definecolor{ipssiBleuDroite}{HTML}{3498DB} % Droite de régression

% ============================================================================
% ENCADRÉS PÉDAGOGIQUES (tcolorbox)
% ============================================================================
\usepackage[most]{tcolorbox}

% Style de base partagé
\tcbset{
    ipssibase/.style={
        enhanced,
        breakable,
        left=4mm,
        right=4mm,
        top=3mm,
        bottom=3mm,
        boxrule=0pt,
        frame hidden,
        borderline west={3pt}{0pt}{#1},
        colback=#1!5,
        colframe=#1,
        fonttitle=\bfseries\sffamily,
        coltitle=#1!80!black,
        attach boxed title to top left={yshift=-2mm, xshift=4mm},
        boxed title style={
            colback=#1!15,
            colframe=#1,
            boxrule=0.5pt,
            arc=2pt,
        },
        arc=2pt,
        outer arc=2pt,
    }
}

% Icônes TikZ pour les encadrés (compatible pdflatex, pas d'emoji Unicode)
\newcommand{\icoDefinition}{%
    \tikz[baseline=-0.3em]{\draw[ipssiBleu, thick] (0,0) -- (0.3,0) -- (0.3,0.35) -- (0,0.35) -- cycle;
    \draw[ipssiBleu, thick] (0.05,0.1) -- (0.25,0.1); \draw[ipssiBleu, thick] (0.05,0.2) -- (0.2,0.2);}%
}
\newcommand{\icoRetenir}{%
    \tikz[baseline=-0.2em]{\draw[ipssiVert, thick, fill=ipssiVert!20] (0.15,0) -- (0.3,0.3) -- (0,0.3) -- cycle;}%
}
\newcommand{\icoAttention}{%
    \tikz[baseline=-0.2em]{\draw[ipssiOrange, thick] (0.15,0.35) -- (0,0) -- (0.3,0) -- cycle;
    \node[ipssiOrange, font=\tiny\bfseries] at (0.15,0.12) {!};}%
}
\newcommand{\icoIntuition}{%
    \tikz[baseline=-0.2em]{\draw[ipssiViolet, thick, fill=ipssiViolet!15]
    (0.15,0.35) -- (0.08,0.2) -- (0.08,0.15) -- (0.22,0.15) -- (0.22,0.2) -- cycle;
    \draw[ipssiViolet, thick] (0.11,0.12) -- (0.11,0.05); \draw[ipssiViolet, thick] (0.19,0.12) -- (0.19,0.05);}%
}
\newcommand{\icoExemple}{%
    \tikz[baseline=-0.2em]{\draw[ipssiCyan, thick] (0,0) rectangle (0.3,0.3);
    \draw[ipssiCyan, thick] (0,0.15) -- (0.3,0.15); \draw[ipssiCyan, thick] (0.15,0) -- (0.15,0.3);}%
}
\newcommand{\icoPython}{%
    \tikz[baseline=-0.2em]{\node[font=\small\bfseries\ttfamily, ipssiVert!70!black] at (0,0) {\textgreater\textgreater};}%
}

% Définition (bleu)
\newtcolorbox{definitionbox}[1][D\'{e}finition]{
    ipssibase=ipssiBleu,
    title={\icoDefinition\; #1},
}

% À retenir (vert)
\newtcolorbox{retenirbox}[1][\`{A} retenir]{
    ipssibase=ipssiVert,
    title={\icoRetenir\; #1},
}

% Attention / Piège (orange)
\newtcolorbox{attentionbox}[1][Attention]{
    ipssibase=ipssiOrange,
    title={\icoAttention\; #1},
}

% Intuition (violet)
\newtcolorbox{intuitionbox}[1][Intuition]{
    ipssibase=ipssiViolet,
    title={\icoIntuition\; #1},
}

% Exemple (cyan)
\newtcolorbox{exemplebox}[1][Exemple]{
    ipssibase=ipssiCyan,
    title={\icoExemple\; #1},
}

% Code Python (style terminal)
\newtcolorbox{pythonbox}[1][Code Python]{
    enhanced,
    breakable,
    left=4mm, right=4mm, top=3mm, bottom=3mm,
    boxrule=0.5pt,
    colback=ipssiGrisClair,
    colframe=ipssiGrisFonce!40,
    fonttitle=\bfseries\ttfamily\small,
    coltitle=ipssiGrisFonce,
    title={\icoPython\; #1},
    attach boxed title to top left={yshift=-2mm, xshift=4mm},
    boxed title style={colback=ipssiGrisClair, colframe=ipssiGrisFonce!40, boxrule=0.5pt, arc=2pt},
    arc=2pt, outer arc=2pt,
}

% ============================================================================
% CODE PYTHON (listings)
% ============================================================================
\usepackage{listings}

\lstdefinestyle{python}{
    language=Python,
    basicstyle=\ttfamily\small,
    keywordstyle=\color{ipssiBleu}\bfseries,
    stringstyle=\color{ipssiVert!70!black},
    commentstyle=\color{gray}\itshape,
    numberstyle=\tiny\color{gray},
    numbers=left,
    numbersep=8pt,
    backgroundcolor=\color{ipssiGrisClair},
    frame=none,
    rulecolor=\color{ipssiGrisFonce!20},
    breaklines=true,
    breakatwhitespace=true,
    tabsize=4,
    showstringspaces=false,
    captionpos=b,
    aboveskip=8pt,
    belowskip=8pt,
    xleftmargin=10pt,
    framexleftmargin=10pt,
    morekeywords={as, with, True, False, None, self, np, plt, import, from},
    literate=
        {é}{{\'e}}1 {è}{{\`e}}1 {ê}{{\^e}}1 {ë}{{\"e}}1
        {à}{{\`a}}1 {â}{{\^a}}1
        {ù}{{\`u}}1 {û}{{\^u}}1
        {ô}{{\^o}}1
        {î}{{\^i}}1 {ï}{{\"i}}1
        {ç}{{\c{c}}}1
        {É}{{\'E}}1 {È}{{\`E}}1
        {→}{{$\rightarrow$}}1
        {≈}{{$\approx$}}1
        {≤}{{$\leq$}}1
        {≥}{{$\geq$}}1
        {★}{{$\bigstar$}}1
        {ŷ}{{$\hat{y}$}}1
        {·}{{$\cdot$}}1,
}

\lstset{style=python}

% Style pour le code inline
\newcommand{\pyinline}[1]{\lstinline[style=python,basicstyle=\ttfamily\small]{#1}}

% ============================================================================
% GRAPHIQUES
% ============================================================================
\usepackage{tikz}
\usepackage{pgfplots}
\pgfplotsset{compat=1.18}
\usetikzlibrary{arrows.meta, calc, positioning, shapes.geometric, decorations.pathreplacing}

% ============================================================================
% TABLEAUX
% ============================================================================
\usepackage{booktabs}
\usepackage{tabularx}
\usepackage{multirow}

% ============================================================================
% LISTES
% ============================================================================
\usepackage{enumitem}
\setlist[itemize]{itemsep=2pt, parsep=0pt, topsep=4pt}
\setlist[enumerate]{itemsep=2pt, parsep=0pt, topsep=4pt}

% ============================================================================
% HYPERLIENS
% ============================================================================
\usepackage[
    colorlinks=true,
    linkcolor=ipssiBleu!80!black,
    urlcolor=ipssiBleu!80!black,
    citecolor=ipssiVert!80!black,
    bookmarks=true,
    bookmarksopen=true,
    pdfauthor={Mohamed EL AFRIT},
    pdftitle={Mathématiques appliquées à l'IA — IPSSI},
    pdfsubject={Formation Bachelor 2 — IPSSI 2025-2026},
]{hyperref}

% ============================================================================
% IMAGES
% ============================================================================
\usepackage{graphicx}

% ============================================================================
% VARIABLES PARAMÉTRABLES (jour et thème)
% ============================================================================
\newcommand{\leJour}{X}
\newcommand{\leTheme}{Thème}
\newcommand{\jourNumero}[1]{\renewcommand{\leJour}{#1}}
\newcommand{\jourTheme}[1]{\renewcommand{\leTheme}{#1}}

% ============================================================================
% EN-TÊTE ET PIED DE PAGE (fancyhdr)
% ============================================================================
\usepackage{fancyhdr}

\pagestyle{fancy}
\fancyhf{}

% En-tête
\fancyhead[L]{\small\sffamily\textcolor{gray}{IPSSI — Maths appliquées à l'IA}}
\fancyhead[R]{\small\sffamily\textcolor{gray}{Jour \leJour\ — \leTheme}}
\renewcommand{\headrulewidth}{0.4pt}

% Pied de page
\fancyfoot[L]{\footnotesize\textcolor{gray}{Mohamed EL AFRIT\enspace\textbullet\enspace\url{www.mohamedelafrit.com}\enspace\textbullet\enspace CC BY-NC-SA 4.0}}
\fancyfoot[C]{\footnotesize\textcolor{gray}{\thepage}}
\fancyfoot[R]{\footnotesize\textcolor{gray}{2025–2026}}
\renewcommand{\footrulewidth}{0.2pt}

% Style pour la page de garde (pas d'en-tête)
\fancypagestyle{pagegarde}{
    \fancyhf{}
    \renewcommand{\headrulewidth}{0pt}
    \fancyfoot[C]{\footnotesize\textcolor{gray}{\thepage}}
    \renewcommand{\footrulewidth}{0pt}
}

% ============================================================================
% PAGE DE GARDE — \pagegarde{titre}{jour}{theme}
% ============================================================================
\newcommand{\pagegarde}[3]{%
    \thispagestyle{pagegarde}
    \begin{center}
        \vspace*{1.5cm}

        % Logo / Bandeau institution
        {\Large\sffamily\bfseries\textcolor{ipssiBleu}{IPSSI}}\par
        \vspace{2pt}
        {\normalsize\sffamily 2\textsuperscript{e} année Bachelor Informatique}\par

        \vspace{0.8cm}
        {\color{ipssiBleu}\rule{0.6\textwidth}{1.5pt}}\par
        \vspace{0.8cm}

        % Titre du module
        {\LARGE\sffamily\bfseries Mathématiques appliquées\\[4pt] à l'Intelligence Artificielle}\par

        \vspace{1cm}

        % Titre du document
        \begin{tcolorbox}[
            enhanced,
            colback=ipssiBleu!5,
            colframe=ipssiBleu!50,
            boxrule=1pt,
            arc=4pt,
            width=0.85\textwidth,
            halign=center,
            top=10pt, bottom=10pt,
        ]
            {\huge\sffamily\bfseries\textcolor{ipssiBleu!80!black}{#1}}\par
            \vspace{6pt}
            {\Large\sffamily Jour #2 — #3}
        \end{tcolorbox}

        \vspace{1.5cm}

        % Informations auteur
        {\large\sffamily
            \textbf{Auteur :} Mohamed EL AFRIT\\[4pt]
            \url{www.mohamedelafrit.com}\\[8pt]
            \textbf{Année académique :} 2025–2026
        }\par

        \vspace{1.5cm}

        % Licence Creative Commons
        {\color{ipssiBleu}\rule{0.4\textwidth}{0.5pt}}\par
        \vspace{8pt}
        {\small\sffamily
            \textcopyright\ 2025 Mohamed EL AFRIT — IPSSI\\[2pt]
            Ce contenu est distribué sous licence\\[2pt]
            \textbf{Creative Commons BY-NC-SA 4.0}\\[2pt]
            \url{https://creativecommons.org/licenses/by-nc-sa/4.0/deed.fr}
        }\par
    \end{center}
    \newpage
}

% ============================================================================
% NIVEAUX D'EXERCICES
% ============================================================================

% Étoile compatible pdflatex
\newcommand{\starfull}[1]{\textcolor{#1}{$\bigstar$}}
\newcommand{\starempty}{\textcolor{gray!40}{$\bigstar$}}

% Fondamental (1 étoile)
\newcommand{\exofondamental}{%
    \starfull{ipssiVert}\starempty\starempty\enspace%
    {\small\sffamily\textcolor{ipssiVert}{\textbf{Fondamental}}}%
}

% Intermédiaire (2 étoiles)
\newcommand{\exointermediaire}{%
    \starfull{ipssiOrange}\starfull{ipssiOrange}\starempty\enspace%
    {\small\sffamily\textcolor{ipssiOrange}{\textbf{Interm\'{e}diaire}}}%
}

% Avancé (3 étoiles)
\newcommand{\exoavance}{%
    \starfull{ipssiRouge}\starfull{ipssiRouge}\starfull{ipssiRouge}\enspace%
    {\small\sffamily\textcolor{ipssiRouge}{\textbf{Avanc\'{e}}}}%
}

% Commande d'exercice numéroté
\newcounter{exocount}[section]
\newcommand{\exercice}[1]{%
    \stepcounter{exocount}%
    \vspace{8pt}%
    \noindent\textbf{\sffamily Exercice \theexocount}\enspace — \enspace #1\par
    \vspace{4pt}%
}

% ============================================================================
% COMMANDES UTILITAIRES
% ============================================================================

% Séparateur visuel entre sections
\newcommand{\separateur}{%
    \vspace{6pt}
    \begin{center}
        {\color{ipssiBleu!30}\rule{0.3\textwidth}{0.5pt}}
    \end{center}
    \vspace{6pt}
}

% Mot-clé en gras coloré
\newcommand{\motcle}[1]{\textbf{\textcolor{ipssiBleu!80!black}{#1}}}

% Formule encadrée importante
\newcommand{\formule}[1]{%
    \begin{center}
        \fcolorbox{ipssiBleu!50}{ipssiBleu!5}{%
            \quad$\displaystyle #1$\quad%
        }
    \end{center}
}

% ============================================================================
% FIN DU PRÉAMBULE
% ============================================================================


\jourNumero{4}
\jourTheme{R\'{e}seaux de neurones}

\begin{document}

\pagegarde{Fiche de synth\`{e}se}{4}{R\'{e}seaux de neurones}

\setstretch{1.05}
\setlength{\parskip}{2pt}
\setlength{\abovedisplayskip}{5pt}
\setlength{\belowdisplayskip}{5pt}

% ======================================================================
\section{Le neurone artificiel}
% ======================================================================

Le neurone artificiel est une version enrichie du perceptron :
il remplace la fonction signe par une \motcle{fonction d'activation} non lin\'{e}aire.

\begin{definitionbox}[Neurone artificiel]
    Un neurone calcule deux \'{e}tapes :
    \begin{align}
        \text{Somme pond\'{e}r\'{e}e :}\quad & z = \vw^T \vx + b
        \label{eq:somme} \\
        \text{Activation :}\quad & a = \sigma(z)
        \label{eq:activation}
    \end{align}
    o\`{u} $\sigma$ est une fonction d'activation non lin\'{e}aire.
\end{definitionbox}

\begin{center}
\begin{tikzpicture}[
    node distance=1.6cm and 2.2cm,
    neuron/.style={circle, draw=ipssiBleu, thick, minimum size=1.1cm, font=\small\sffamily},
    input/.style={circle, draw=gray, thick, minimum size=0.8cm, font=\small},
    arr/.style={->, thick, >=stealth},
]
    % Entrees
    \node[input] (x1) at (0, 1.2) {$x_1$};
    \node[input] (x2) at (0, 0) {$x_2$};
    \node[input] (xp) at (0, -1.2) {$x_p$};
    \node[gray, font=\small] at (0, -0.6) {$\vdots$};

    % Neurone
    \node[neuron, fill=ipssiBleu!10] (n) at (3.2, 0) {$\Sigma$};

    % Activation
    \node[neuron, fill=ipssiViolet!15, draw=ipssiViolet] (act) at (5.8, 0) {$\sigma$};

    % Sortie
    \node[font=\small\sffamily] (out) at (7.8, 0) {$a = \sigma(z)$};

    % Fleches
    \draw[arr] (x1) -- node[above, font=\scriptsize, pos=0.4] {$w_1$} (n);
    \draw[arr] (x2) -- node[above, font=\scriptsize, pos=0.4] {$w_2$} (n);
    \draw[arr] (xp) -- node[below, font=\scriptsize, pos=0.4] {$w_p$} (n);
    \draw[arr] (n) -- node[above, font=\scriptsize] {$z$} (act);
    \draw[arr] (act) -- (out);

    % Biais
    \node[font=\scriptsize, gray] (bias) at (3.2, -1.5) {biais $b$};
    \draw[arr, gray, dashed] (bias) -- (n);

    % Labels
    \node[below, font=\scriptsize\sffamily, gray] at (1.5, -2) {$z = w_1 x_1 + w_2 x_2 + \cdots + w_p x_p + b$};
\end{tikzpicture}
\end{center}

% ======================================================================
\section{Fonctions d'activation}
% ======================================================================

\begin{center}
\renewcommand{\arraystretch}{1.4}
\begin{tabularx}{\textwidth}{l >{\centering\arraybackslash}p{4.5cm} X}
    \toprule
    \textbf{Fonction} & \textbf{Formule} & \textbf{Propri\'{e}t\'{e}s} \\
    \midrule
    Sigmo\"{\i}de
        & $\sigma(z) = \dfrac{1}{1 + e^{-z}}$
        & Sortie dans $(0,1)$. Id\'{e}ale pour les probabilit\'{e}s.
          D\'{e}riv\'{e}e : $\sigma'(z) = \sigma(z)(1-\sigma(z))$. \\[6pt]
    ReLU
        & $\text{ReLU}(z) = \max(0, z)$
        & Sortie $\geq 0$. Rapide \`{a} calculer. La plus utilis\'{e}e
          dans les r\'{e}seaux modernes. \\[6pt]
    Signe
        & $\text{signe}(z) \in \{0, 1\}$
        & \'{E}chelon binaire (perceptron Jour~3).
          Non diff\'{e}rentiable $\Rightarrow$ pas de gradient. \\
    \bottomrule
\end{tabularx}
\end{center}

\begin{center}
\begin{tikzpicture}[scale=0.72]
    % --- Sigmoide ---
    \begin{scope}[xshift=0cm]
        \draw[->, thick] (-3.5,0) -- (3.8,0) node[right, font=\scriptsize] {$z$};
        \draw[->, thick] (0,-0.3) -- (0,2.5) node[above, font=\scriptsize] {$\sigma(z)$};
        \draw[ipssiBleu, very thick, domain=-3.3:3.3, samples=60]
            plot (\x, {2/(1+exp(-\x))});
        \draw[gray, dashed, thin] (-3.5,2) -- (3.5,2);
        \draw[gray, dashed, thin] (-3.5,1) -- (3.5,1);
        \node[font=\scriptsize, gray] at (3, 2.2) {$1$};
        \node[font=\scriptsize, gray] at (3, 1.2) {$0{,}5$};
        \node[below, font=\small\sffamily\bfseries, ipssiBleu] at (0, -0.7) {Sigmo\"{\i}de};
    \end{scope}
    % --- ReLU ---
    \begin{scope}[xshift=9cm]
        \draw[->, thick] (-3.5,0) -- (3.8,0) node[right, font=\scriptsize] {$z$};
        \draw[->, thick] (0,-0.3) -- (0,3.5) node[above, font=\scriptsize] {ReLU};
        \draw[ipssiRouge, very thick] (-3.3,0) -- (0,0);
        \draw[ipssiRouge, very thick] (0,0) -- (3,3);
        \node[below, font=\small\sffamily\bfseries, ipssiRouge] at (0, -0.7) {ReLU};
    \end{scope}
\end{tikzpicture}
\end{center}

\begin{retenirbox}[Pourquoi les activations sont indispensables]
    Sans fonction d'activation, un empilement de couches lin\'{e}aires reste \textbf{lin\'{e}aire}.
    Les activations introduisent la \textbf{non-lin\'{e}arit\'{e}} qui permet au r\'{e}seau de
    mod\'{e}liser des relations complexes (comme le XOR).
\end{retenirbox}

% ======================================================================
\section{R\'{e}seau multicouche (MLP)}
% ======================================================================

\begin{definitionbox}[Perceptron multicouche]
    Un \textbf{MLP} (\emph{Multi-Layer Perceptron}) est compos\'{e} de :
    \begin{itemize}
        \item \textbf{Couche d'entr\'{e}e} : re\c{c}oit les features $\vx$ (pas de calcul)
        \item \textbf{Couche(s) cach\'{e}e(s)} : neurones qui cr\'{e}ent de nouvelles repr\'{e}sentations
        \item \textbf{Couche de sortie} : produit la pr\'{e}diction $\yhat$
    \end{itemize}
    Chaque neurone est connect\'{e} \`{a} \textbf{tous} les neurones de la couche suivante.
\end{definitionbox}

\begin{center}
\begin{tikzpicture}[
    node distance=1.2cm and 2.5cm,
    neuron/.style={circle, draw, thick, minimum size=0.9cm, font=\scriptsize\sffamily},
    arr/.style={->, thick, >=stealth, gray!70},
    scale=0.85, transform shape,
]
    % Entrees
    \node[neuron, draw=gray, fill=gray!10] (i1) at (0, 1) {$x_1$};
    \node[neuron, draw=gray, fill=gray!10] (i2) at (0, -1) {$x_2$};
    \node[below, font=\scriptsize\sffamily, gray] at (0, -2) {Entr\'{e}e};

    % Couche cachee
    \node[neuron, draw=ipssiBleu, fill=ipssiBleu!10] (h1) at (3, 1.5) {$h_1$};
    \node[neuron, draw=ipssiBleu, fill=ipssiBleu!10] (h2) at (3, 0) {$h_2$};
    \node[neuron, draw=ipssiBleu, fill=ipssiBleu!10] (h3) at (3, -1.5) {$h_3$};
    \node[below, font=\scriptsize\sffamily, ipssiBleu] at (3, -2.3) {Cach\'{e}e};

    % Sortie
    \node[neuron, draw=ipssiVert, fill=ipssiVert!10] (o1) at (6, 0) {$\yhat$};
    \node[below, font=\scriptsize\sffamily, ipssiVert] at (6, -2) {Sortie};

    % Connexions entree -> cachee
    \foreach \i in {1,2} \foreach \h in {1,2,3}
        \draw[arr] (i\i) -- (h\h);
    % Connexions cachee -> sortie
    \foreach \h in {1,2,3}
        \draw[arr] (h\h) -- (o1);

    % Labels
    \node[above, font=\scriptsize, ipssiBleu] at (1.5, 1.7) {$\mathbf{W}^{(1)}, \mathbf{b}^{(1)}$};
    \node[above, font=\scriptsize, ipssiVert] at (4.5, 1.2) {$\mathbf{w}^{(2)}, b^{(2)}$};
\end{tikzpicture}
\end{center}

\subsection{Forward pass (propagation avant)}

Le calcul se fait \textbf{couche par couche}, de l'entr\'{e}e vers la sortie :

\begin{align}
    \text{Couche cach\'{e}e :}\quad
        & \mathbf{h} = \sigma\bigl(\mathbf{W}^{(1)} \vx + \mathbf{b}^{(1)}\bigr)
        \label{eq:forward1} \\[4pt]
    \text{Couche de sortie :}\quad
        & \yhat = \sigma\bigl(\mathbf{w}^{(2)T} \mathbf{h} + b^{(2)}\bigr)
        \label{eq:forward2}
\end{align}

\begin{exemplebox}[R\'{e}seau 2-2-1 --- calcul complet]
    Entr\'{e}e $\vx = (5, 4)^T$ :

    \smallskip
    $z_1 = 0{,}3 \times 5 + 0{,}5 \times 4 - 2 = 1{,}5 \;\Rightarrow\; h_1 = \sigma(1{,}5) = 0{,}82$

    $z_2 = {-}0{,}2 \times 5 + 0{,}4 \times 4 + 0{,}5 = 1{,}1 \;\Rightarrow\; h_2 = \sigma(1{,}1) = 0{,}75$

    $z_{\text{out}} = 0{,}8 \times 0{,}82 + 0{,}6 \times 0{,}75 - 0{,}5 = 0{,}61 \;\Rightarrow\; \yhat = \sigma(0{,}61) = 0{,}65$

    \smallskip
    Pr\'{e}diction : 65\% de probabilit\'{e} d'\^{e}tre en classe~1.
\end{exemplebox}

% ======================================================================
\section{Apprentissage et backpropagation}
% ======================================================================

\subsection{Fonctions de perte}

\begin{center}
\renewcommand{\arraystretch}{1.35}
\begin{tabularx}{\textwidth}{l >{\centering\arraybackslash}p{5.5cm} X}
    \toprule
    \textbf{T\^{a}che} & \textbf{Perte} & \textbf{Usage} \\
    \midrule
    R\'{e}gression
        & $J = \frac{1}{n}\sum(y_i - \yhat_i)^2$
        & MSE (Jours 1--2) \\
    Classification
        & $L = -\frac{1}{n}\sum\bigl[y_i\ln\yhat_i + (1{-}y_i)\ln(1{-}\yhat_i)\bigr]$
        & Cross-entropy binaire \\
    \bottomrule
\end{tabularx}
\end{center}

\subsection{Backpropagation (intuition)}

\begin{intuitionbox}[Propager l'erreur en arri\`{e}re]
    La \textbf{backpropagation} calcule la contribution de \textbf{chaque poids} \`{a} l'erreur
    finale, en remontant de la sortie vers l'entr\'{e}e. C'est la \textbf{r\`{e}gle de la cha\^{i}ne}
    (Jour~2) appliqu\'{e}e \`{a} chaque couche du r\'{e}seau.
\end{intuitionbox}

\begin{retenirbox}[La boucle d'entra\^{i}nement]
    L'entra\^{i}nement d'un r\'{e}seau r\'{e}p\`{e}te 3 \'{e}tapes :
    \begin{enumerate}
        \item \textbf{Forward pass} : calculer $\yhat$ couche par couche
        \item \textbf{Backward pass} : calculer $\nabla L$ pour chaque poids (backpropagation)
        \item \textbf{Mise \`{a} jour} : $\boldsymbol{\theta} \leftarrow \boldsymbol{\theta} - \alpha \nabla L$
    \end{enumerate}
    C'est la \textbf{m\^{e}me boucle} qu'au Jour~2, mais appliqu\'{e}e \`{a} toutes les couches.
\end{retenirbox}

% ======================================================================
\section{Comparatif : perceptron vs r\'{e}seau}
% ======================================================================

\begin{center}
\renewcommand{\arraystretch}{1.4}
\begin{tabularx}{\textwidth}{l >{\centering\arraybackslash}X >{\centering\arraybackslash}X}
    \toprule
    & \textbf{Perceptron (Jour 3)} & \textbf{R\'{e}seau MLP (Jour 4)} \\
    \midrule
    Architecture  & 1 couche (entr\'{e}e $\to$ sortie) & Plusieurs couches \\
    Activation    & Signe (0 ou 1) & Sigmo\"{\i}de, ReLU, \ldots \\
    Fronti\`{e}re & Lin\'{e}aire (droite) & \textbf{Non lin\'{e}aire} (courbe) \\
    R\'{e}sout XOR ? & Non & \textbf{Oui} \\
    Accuracy dataset & ${\sim}87$--$90\%$ & ${\sim}93$--$97\%$ \\
    Apprentissage & Correction d'erreurs & Backpropagation + GD \\
    scikit-learn  & \texttt{Perceptron} & \texttt{MLPClassifier} \\
    \bottomrule
\end{tabularx}
\end{center}

% ======================================================================
\section{Limites et risques}
% ======================================================================

\begin{attentionbox}[Surapprentissage (overfitting)]
    Un r\'{e}seau trop gros peut \textbf{m\'{e}moriser} les donn\'{e}es au lieu de \textbf{g\'{e}n\'{e}raliser}.
    Solutions : jeu de \textbf{validation}, \textbf{r\'{e}gularisation} ($L_2$), \textbf{dropout},
    \textbf{early stopping}.
\end{attentionbox}

\begin{attentionbox}[Biais dans les donn\'{e}es]
    Un r\'{e}seau ne fait qu'apprendre les \textbf{patterns pr\'{e}sents dans les donn\'{e}es}.
    Si les donn\'{e}es contiennent des biais (discrimination), le mod\`{e}le les reproduira.
    La responsabilit\'{e} \'{e}thique incombe aux humains qui con\c{c}oivent le syst\`{e}me.
\end{attentionbox}

% ======================================================================
\section{Formules essentielles \`{a} retenir}
% ======================================================================

\begin{tcolorbox}[
    enhanced, colback=ipssiBleu!3, colframe=ipssiBleu!40, boxrule=1pt, arc=4pt,
    title={\sffamily\bfseries Formules du Jour 4}, coltitle=white, colbacktitle=ipssiBleu!70,
]
\begin{align}
    \text{Neurone}\quad
        & a = \sigma(\vw^T \vx + b)
        \label{eq:f-neurone} \\[4pt]
    \text{Sigmo\"{\i}de}\quad
        & \sigma(z) = \frac{1}{1 + e^{-z}}, \quad \sigma'(z) = \sigma(z)(1{-}\sigma(z))
        \label{eq:f-sigmoid} \\[4pt]
    \text{ReLU}\quad
        & \text{ReLU}(z) = \max(0, z)
        \label{eq:f-relu} \\[4pt]
    \text{Forward (cach\'{e}e)}\quad
        & \mathbf{h} = \sigma(\mathbf{W}^{(1)} \vx + \mathbf{b}^{(1)})
        \label{eq:f-fwd1} \\[4pt]
    \text{Forward (sortie)}\quad
        & \yhat = \sigma(\mathbf{w}^{(2)T} \mathbf{h} + b^{(2)})
        \label{eq:f-fwd2} \\[4pt]
    \text{Cross-entropy}\quad
        & L = -\tfrac{1}{n}\sum\bigl[y_i\ln\yhat_i + (1{-}y_i)\ln(1{-}\yhat_i)\bigr]
        \label{eq:f-ce} \\[4pt]
    \text{Mise \`{a} jour}\quad
        & \boldsymbol{\theta} \leftarrow \boldsymbol{\theta} - \alpha \nabla L
        \label{eq:f-update}
\end{align}
\end{tcolorbox}

\separateur

\begin{center}
    \small\sffamily
    \textbf{Pipeline du Jour 4 :}\quad
    \fcolorbox{ipssiBleu!40}{ipssiBleu!5}{\;$\vx$\;}
    $\;\xrightarrow{\mathbf{W}^{(1)}}\;$
    \fcolorbox{ipssiBleu!40}{ipssiBleu!5}{\;$\sigma$\;}
    $\;\to\;$
    \fcolorbox{ipssiViolet!40}{ipssiViolet!5}{\;$\mathbf{h}$\;}
    $\;\xrightarrow{\mathbf{w}^{(2)}}\;$
    \fcolorbox{ipssiBleu!40}{ipssiBleu!5}{\;$\sigma$\;}
    $\;\to\;$
    \fcolorbox{ipssiVert!40}{ipssiVert!5}{\;$\yhat$\;}
\end{center}

\vspace{6pt}

\begin{center}
    \small\sffamily
    \textbf{Les 4 jours en un mot :}\quad
    \fcolorbox{gray!30}{gray!5}{\;J1 Repr\'{e}sentation\;}
    $\;\to\;$
    \fcolorbox{gray!30}{gray!5}{\;J2 Optimisation\;}
    $\;\to\;$
    \fcolorbox{gray!30}{gray!5}{\;J3 Classification\;}
    $\;\to\;$
    \fcolorbox{gray!30}{gray!5}{\;J4 R\'{e}seaux\;}
\end{center}

\vfill

\begin{center}
    {\footnotesize\sffamily\textcolor{gray}{%
        \textcopyright\ 2025 Mohamed EL AFRIT --- IPSSI \enspace\textbullet\enspace
        \url{www.mohamedelafrit.com} \enspace\textbullet\enspace
        CC BY-NC-SA 4.0
    }}
\end{center}

\end{document}
